\chapter{Especificação da MiniLang}

\section{Vocabulário geral}

O alfabeto geral da MiniLang utiliza letras ASCII, dígitos, sublinhado, operadores, delimitadores, ponto decimal, aspas e espaços de acordo com o papel de cada categoria lexical. Palavras reservadas são reconhecidas por uma tabela auxiliar e não constituem uma das seis ERs principais.

\section{Palavras reservadas}

A tabela de palavras reservadas contém:

\begin{center}
\texttt{program}, \texttt{int}, \texttt{float}, \texttt{string}, \texttt{bool}, \texttt{time},

\texttt{if}, \texttt{else}, \texttt{while}, \texttt{print}, \texttt{true}, \texttt{false}.
\end{center}

\section{Tokens principais}

Os tokens definidos pelo projeto incluem \texttt{PROGRAM}, tipos, estruturas de controle, \texttt{BOOLEAN}, identificadores, inteiros, decimais, notação científica, strings, horários, operadores e delimitadores. \texttt{WHITESPACE} e \texttt{COMMENT} são classes internas e são descartadas da saída final.

\section{Seis Expressões Regulares principais}

As seis ERs principais são:

\begin{center}
\small
\begin{tabular}{|c|r|r|r|r|}
\toprule
ER & Estados & Transições & $\varepsilon$ & Rotuladas \\
\midrule
ER-01 & 28 & 35 & 27 & 8 \\
ER-02 & 10 & 12 & 9 & 3 \\
ER-03 & 18 & 22 & 16 & 6 \\
ER-04 & 42 & 52 & 40 & 12 \\
ER-05 & 8 & 9 & 6 & 3 \\
ER-06 & 16 & 16 & 9 & 7 \\
\bottomrule
\end{tabular}
\end{center}

As contagens acima são as contagens definitivas congeladas no projeto.

\section{Ordem de tokenização}

A ordem oficial é:

\begin{enumerate}
    \item whitespace;
    \item comentário;
    \item string;
    \item operador composto;
    \item número científico;
    \item número decimal;
    \item horário;
    \item número inteiro;
    \item identificador ou palavra reservada;
    \item operador simples;
    \item delimitador;
    \item erro léxico.
\end{enumerate}

A prioridade é combinada com o princípio de maior lexema. Em especial, o lexer não deve transformar um candidato numérico inválido em vários tokens menores.

\section{Regras auxiliares importantes}

Comentários iniciados por \texttt{//} são ignorados antes que \texttt{/} seja tratado como divisão. Operadores compostos \texttt{==}, \texttt{!=}, \texttt{<=} e \texttt{>=} têm prioridade sobre operadores simples. A string é analisada antes das regras internas que poderiam interpretar seus símbolos como operadores. O ponto isolado não é um token válido.
